\par\medskip
\hypertarget{native-073-article055}{}\pdfbookmark[2]{55.}{article-055}
\begin{center}55.\end{center}
As above, let us denote the longitude of the ascending node of the orbit on the ecliptic by $\AscendingNode$ and the inclination by $i$. Let the longitude of the ascending node of the new plane on the ecliptic be denoted by $n$ and its inclination by $\epsilon$. Let the distance of the ascending node of the orbit on the new plane from the ascending node of the new plane on the ecliptic (the arc $n\AscendingNode'$ in Fig. 2) be denoted by $\AscendingNode'$, and the inclination of the orbit to the new plane by $i'$. Finally, let the arc from $\AscendingNode$ to $\AscendingNode'$ in the direction of motion be denoted by $\Delta$. Our spherical triangle then has sides $\AscendingNode-n$, $\AscendingNode'$, and $\Delta$, with opposite angles $i'$, $180^0-i$, and $\epsilon$. The formulas of the preceding article then give
\SourceLeaf{074}
\[
\sin\tfrac{1}{2}i'\sin\tfrac{1}{2}(\AscendingNode'+\Delta) = \sin\tfrac{1}{2}(\AscendingNode-n)\sin\tfrac{1}{2}(i+\epsilon)
\]
\[
\sin\tfrac{1}{2}i'\cos\tfrac{1}{2}(\AscendingNode'+\Delta) = \cos\tfrac{1}{2}(\AscendingNode-n)\sin\tfrac{1}{2}(i-\epsilon)
\]
\[
\cos\tfrac{1}{2}i'\sin\tfrac{1}{2}(\AscendingNode'-\Delta) = \sin\tfrac{1}{2}(\AscendingNode-n)\cos\tfrac{1}{2}(i+\epsilon)
\]
\[
\cos\tfrac{1}{2}i'\cos\tfrac{1}{2}(\AscendingNode'-\Delta) = \cos\tfrac{1}{2}(\AscendingNode-n)\cos\tfrac{1}{2}(i-\epsilon).
\]
The first two equations give $\frac{1}{2}(\AscendingNode'+\Delta)$ and
$\sin\frac{1}{2}i'$. The remaining two give $\frac{1}{2}(\AscendingNode'-\Delta)$ and $\cos\frac{1}{2}i'$.
From $\frac{1}{2}(\AscendingNode'+\Delta)$ and $\frac{1}{2}(\AscendingNode'-\Delta)$ we obtain $\AscendingNode'$ and
$\Delta$. From $\sin\frac{1}{2}i'$ or $\cos\frac{1}{2}i'$ we can determine
$i'$ (agreement between the two determinations provides a check on the calculation). The ambiguity of whether to take $\frac{1}{2}(\AscendingNode'+\Delta)$ and
$\frac{1}{2}(\AscendingNode'-\Delta)$ between $0$ and $180^0$ or between $180^0$ and $360^0$
is resolved by choosing them so that both $\sin\frac{1}{2}i'$ and $\cos\frac{1}{2}i'$ are positive,
since by the nature of the problem $i'$ must be less than $180^0$.
